% Models section (Part D, Milestone 5)
% Include from reports/midsem/main.tex.
%
% Describes the models in src/vio_drift/models/ml/ and the evaluation
% protocol in tuning.py / splits.json. Hyperparameter grids are the
% PARAM_GRID constants in baselines.py, minibatch_gd.py and random_forest.py; chosen
% values are in results/metrics/err_mag_m_tuning.csv.

\section{Models}
\label{sec:models}

Every model maps a window's feature vector $\mathbf{x} \in \mathbb{R}^{104}$
(Section~\ref{sec:feature-engineering}) to a predicted drift $\hat{y}$, the
1\,s position error \texttt{err\_mag\_m} in metres. We compare a mean baseline,
four linear models (OLS, Ridge, Lasso, and a linear model trained by
mini-batch gradient descent) and one non-linear model, a Random Forest. All
hyperparameters are chosen by grid search (Section~\ref{sec:protocol}).

\subsection{Baseline and linear models}

\paragraph{Mean baseline.}
Predict the average training label for every window,
$\hat{y} = \bar{y}_{\text{train}}$. It ignores the features entirely, so any
model that cannot beat it has learned nothing that transfers to a new flight.

\paragraph{Linear regression (OLS).}
Ordinary least squares: the linear model $\hat{y} = \mathbf{w}^\top \mathbf{x} + b$
that minimises the squared training error
$\sum_{i} (y_i - \mathbf{w}^\top \mathbf{x}_i - b)^2$, with nothing to tune.
It is the unregularised reference for Ridge: comparing the two measures what
Ridge's penalty buys rather than assuming it
(Section~\ref{sec:overfitting}).

\paragraph{Ridge regression.}
The same linear model, fitted by minimising
\begin{equation}
  \sum_{i} \bigl(y_i - \mathbf{w}^\top \mathbf{x}_i - b\bigr)^2
  + \alpha \lVert \mathbf{w} \rVert^2 .
\end{equation}
The penalty $\alpha \lVert \mathbf{w} \rVert^2$ shrinks the weights. It is
needed here because many features are near-duplicates (e.g.\ variance is the
square of standard deviation; 17 feature pairs have $|r| > 0.95$): without it,
a least-squares fit would give such pairs large weights of opposite sign that
cancel on the training data and blow up on new data (as OLS does,
Section~\ref{sec:overfitting}). Features are standardised
(zero mean, unit variance) first, so the penalty treats every feature equally.
We tune $\alpha \in \{10^{-2}, 10^{-1}, \dots, 10^{5}\}$ (8 values).

\paragraph{Lasso regression.}
The same model with an L1 penalty instead,
\begin{equation}
  \frac{1}{2n}\sum_{i} \bigl(y_i - \mathbf{w}^\top \mathbf{x}_i - b\bigr)^2
  + \alpha \lVert \mathbf{w} \rVert_1 ,
  \qquad \lVert \mathbf{w} \rVert_1 = \textstyle\sum_j |w_j| .
\end{equation}
L2 shrinks every weight a little; L1 pushes weak weights to \emph{exactly} zero,
so Lasso also selects features. With redundant features this is a natural fit:
from a group of near-duplicates it tends to keep one. The useful range of
$\alpha$ is small because the label's spread is only about $0.015$\,m; we tune
$\alpha \in \{10^{-6}, 10^{-5}, \dots, 10^{-1}\}$ (6 values).

\subsection{Mini-batch gradient descent}

OLS, Ridge and Lasso are fitted by a solver. As a fourth linear model we fit
the same $\hat{y} = \mathbf{w}^\top \mathbf{x} + b$ iteratively, with our own
implementation of mini-batch gradient descent. Each epoch shuffles the training
windows and splits them into batches $B$; for each batch the weights take one
step against the gradient of the batch loss
\begin{equation}
  L_B(\mathbf{w}, b) = \frac{1}{2|B|} \sum_{i \in B}
    \bigl(y_i - \mathbf{w}^\top \mathbf{x}_i - b\bigr)^2 + \alpha R(\mathbf{w}),
  \qquad
  \mathbf{w} \leftarrow \mathbf{w} - \eta\, \nabla_{\mathbf{w}} L_B ,
\end{equation}
where $R(\mathbf{w}) = \tfrac12\lVert \mathbf{w} \rVert^2$ (L2) or
$\lVert \mathbf{w} \rVert_1$ (L1, using the subgradient
$\operatorname{sign}(\mathbf{w})$), and $\eta$ is the learning rate. Batch size
sits between stochastic gradient descent (one window per step) and full-batch
gradient descent (all windows per step). Averaging the gradient over a batch
smooths out single unusual windows while still taking many steps per epoch;
notebook~05 shows that at the same learning rate, batch size 1 diverges on our
data (some standardised features have extreme outliers) and the full batch
barely moves in 50 epochs, while batches of 32 converge. Features are
standardised first, and weights start at zero with $b$ at the training mean.
We run 50 epochs and tune $\eta \in \{0.001, 0.01\}$, batch size
$\in \{32, 128\}$, the penalty (L2 or L1) and
$\alpha \in \{10^{-4}, \dots, 10^{-1}\}$: 32 combinations.

\subsection{Random Forest}

A \emph{regression tree} recursively splits the data with yes/no questions on
one feature at a time (``is \texttt{ax\_rms} $> 9.6$?''). At each node it picks
the feature and threshold that most reduce the squared error of the label in
the two resulting groups, and a leaf predicts the mean label of the training
windows that reach it. A single deep tree fits noise. A \emph{Random Forest}
\cite{breiman2001random} averages $B$ trees,
\begin{equation}
  \hat{y}(\mathbf{x}) = \frac{1}{B} \sum_{b=1}^{B} T_b(\mathbf{x}),
\end{equation}
where each tree $T_b$ is trained on a bootstrap resample of the windows and
considers only a random subset of features at each split. These two sources of
randomness make the trees' errors less correlated, so averaging them cuts
variance while keeping each tree's low bias.

Trees depend only on the \emph{order} of a feature's values, so no scaling is
needed. We use $B = 300$ trees and tune the maximum depth
($\{\text{none}, 4, 8\}$), the minimum windows per leaf ($\{1, 5, 20\}$) and the
number of features tried per split (all $p = 104$, $\sqrt{p}$, or $0.3p$): 27
combinations. Shallow trees with large leaves are the forest's main defence
against over-fitting on our small data set.

\subsection{Evaluation protocol}
\label{sec:protocol}

Windows overlap by 75\% (1\,s length, 0.25\,s stride), so neighbouring windows
are nearly identical. A random train/test split would therefore score a model
on near-copies of its training windows. We split by \emph{sequence}
instead, using \emph{nested} leave-one-sequence-out (LOSO) cross-validation:
\begin{itemize}
  \item \textbf{Outer loop} (\texttt{splits.json}): three folds, each training on
        two sequences and testing on the third. Only these held-out scores
        are reported.
  \item \textbf{Inner loop} (\texttt{tuning.py}): within a fold's two training
        sequences, each hyperparameter combination is trained on one sequence
        and scored on the other, both ways. The best combination by mean
        absolute error is refit on both and then scored once on the outer
        test sequence.
\end{itemize}
Every scaler sits inside its model's pipeline, so it is fit on training
windows only. Notebook~05 verifies this for each fold: each saved scaler's mean
equals the training-sequence mean, not the pooled mean. As a further control,
a Random Forest trained on \emph{shuffled} labels does no better than the mean
baseline on held-out data.

We report mean absolute error (MAE), root mean squared error (RMSE), and
$R^2 = 1 - \sum_i (y_i - \hat{y}_i)^2 / \sum_i (y_i - \bar{y})^2$ computed on the
held-out sequence. $R^2 = 0$ means no better than knowing the
\emph{test} sequence's mean. The mean baseline uses the \emph{training} mean,
so its $R^2$ is slightly negative.
